\section{A Hierarchy of Cues}
\label{sec:tiers}

A benchmark can be solvable at several levels, and those levels deserve different verdicts. We fix
a four-tier hierarchy and use it consistently, including when judging our own results.

\begin{description}[noitemsep,topsep=2pt,leftmargin=1.2em]
  \item[Tier 1: variable identity.] \emph{Which} variable symbols appear. Purely notational:
    under a consistent relabelling the expression denotes the same function, so any identification
    that depends on it is a shortcut, always.
  \item[Tier 2: operator and arity inventory.] \emph{Which} operators appear, how many, and over
    how many variables. Semantically meaningful, $x_0/x_1$ really is a different function from
    $x_0 \cdot x_1$, but still bag-level: a model reading only the multiset has not parsed
    anything.
  \item[Tier 3: order, depth and local shape.] Token adjacency, depth profile, parent--child
    operator statistics, subtree sizes. $x_0 + x_1 \cdot x_2$ and $(x_0 + x_1) \cdot x_2$ differ in
    all of these. Crucially, \emph{none of it requires composition}: these are surface statistics of
    a tree, computable without recursion.
  \item[Tier 4: compositional structure.] How the operators actually combine.
\end{description}

\paragraph{Why tier 3 has to exist.} An earlier version of this hierarchy went straight from tier 2
to composition, licensing a bad inference: \emph{encoder beats a bag $\Rightarrow$ encoder learned
structure}. Everything in tier 3 is available to a model that never composes anything, so clearing a
bag establishes only that the representation uses information beyond the tested bag cues.
\S\ref{sec:tier3} measures that gap.

\paragraph{Non-learned is not structure-free.} Our \emph{bag classifiers} (variable-only,
variable-blind, full-token, and the tier-3 baggers) are nearest-centroid over a multiset with no
parameters. The \emph{untrained encoder} is the trained architecture, randomly initialised. Both are
\emph{non-learned}, but only the bags are structure-free: a random Tree-LSTM still has recursive
composition, parent/child interaction and depth sensitivity, computing over structure without having
learned any. Hence \textbf{non-learned baselines}, not ``structure-free'' ones, a distinction
\S\ref{sec:tier3} measures rather than stipulates. Claims must accordingly name their tier: we report
identification \emph{above the strongest non-learned baseline}, and reserve claims about composition
for the twin probe.

\subsection{Setup}
\label{sec:setup}

Expressions are binary trees over $\{+, -, *, /, \sin, \cos, \exp, \log, \sqrt{\ }\}$. The physics
suite is the 10 AI~Feynman equations \citep{udrescu2020ai} with 1--3 variables, each mapped to a
distinct variable subset of a 16-symbol pool; that mapping is the convention under audit.
Equivalence classes come from a 13-schema paraphrase library and have ${\approx}3$ members differing
by a single rewrite, the easiest possible positive, which bounds every number we report. Encoders
are Tree-LSTMs \citep{tai2015improved}, GINs \citep{xu2019powerful} and Transformers
\citep{vaswani2017attention}, 128-dimensional, trained with satisfaction classification plus a
semi-hard triplet loss \citep{schroff2015facenet}. The metric is the fraction of held-out embeddings
whose nearest training-class centroid is correct; chance $1/K$ is a uniform-class conceptual
baseline, not the empirical expectation under every retrieval construction. Every baseline is
evaluated under the identical centroid protocol, which is what makes it comparable to the trained
number rather than merely adjacent to it.

\paragraph{Neither baseline substitutes for the other.} On EQNET \textsc{oneVarPoly13} the bag
scores $0.15$ where the untrained encoder scores $0.66$; on Feynman-10 the ordering reverses ($1.00$
versus $0.13$). A recursive nonlinearity is sensitive to shape a bag discards, and blind to counts a
bag sees, so both are reported, always.
